% -----------------------------------------------------------------------------
% This file is automatically generated.
% Do not edit manually.
% Source: notebooks/support/population-estimation/support.Rmd
% -----------------------------------------------------------------------------


\noindent The point estimate of the mean ($\bar{y}$) is obtained by dividing the sum of $y$ in the sample by the sample size. In \ttblue{R} we simply use the \ttblue{mean} function.

\par\addvspace{\topsep}
\begingroup
\fvset{listparameters={%
  \setlength{\topsep}{0pt}%
  \setlength{\partopsep}{0pt}%
  \setlength{\parsep}{0pt}%
  \setlength{\itemsep}{0pt}%
}}
\begin{Verbatim}[breaklines=true,formatcom=\color{ada_blue}]
(y_bar <- mean(sample1$gross))
\end{Verbatim}
\begin{Verbatim}[breaklines=true,formatcom=\color{ada_light_blue}]
[1] 3014.16
\end{Verbatim}
\endgroup
\par\addvspace{\topsep}
\noindent We can use the sample mean as an estimate for the total payroll in the population. The estimate of the total monthly payroll is then obtained by multiplying the sample mean with the number of elements in the population $N$.
